\documentclass{article}
\usepackage{graphicx} % Required for inserting images
\usepackage[margin=1.2in]{geometry} % Adjust margins
\usepackage{amsmath}
\usepackage{float}

\title{Mocks Revision}
\author{Reuben Okotie}
\date{December 2025}

\setlength{\parindent}{0pt}
\setlength{\parskip}{0.35em}

\begin{document}

\maketitle
I hope for this document to be extensive and include all the subtopics for the A-Level Physics spec. I will vet over this to have accurate notes throughout the course as a whole and hope to use this for the actual A-Level (I might also do this for select maths topics that are quite conceptually difficult like \textit{Vectors}, \textit{Differential Equations} \& \textit{Mechanics}.

% Notes are either with "\textit{}", "\textbf{}" and just notes.

\section{Waves}
\subsection{Progressive Waves}
Progressive waves are waves that transfer energy without transferring material.

\subsubsection{Longitudinal \& Transverse Waves}
The oscillations of transverse waves are perpendicular to the direction of energy transfer. 

The oscillations of longitudinal waves are parallel to the direction of energy transfer. 

\subsubsection{Polarisation}
** Polarisation is such that the plane of wave oscillation only occurs on one plane.

Only transverse waves can be polarised since they exist along a plane (whereas longitudinal waves could only exist along a given line and thus couldn't be `polarised`).

Examples of polarisation: 
\begin{itemize}
    \item \textit{Review}: Polarised glasses have a polarising filter that means that only light that is incident on the given plane that the glasses work for will pass through to the user. \textit{Understand why this can allow people to see through the bottom of lakes}
    \item \textit{Review}: For old TV antennas, the polarisation of the antennas meant that only certain frequencies of the radio waves could be incident (i.e. the radio waves in the desired frequency appeared louder and clearer (\textit{why?}) and allowed for targeted communications.
\end{itemize}

\subsection{Stationary Waves}
Stationary waves are waves that do not transfer energy (\textit{...so then what are they? I can only describe their negation but that isn't really anything})

They are formed by two progressive waves moving in opposite directions interfering both constructively and destructively with the same frequency ($f$), wavelength ($\lambda$) and constant phase difference($\phi$). 

Nodes are points of total destructive interference (on stationary waves) and are points of zero displacement. Anti-nodes are points of constructive interference and are at maximum displacement (or amplitudes of the stationary wave). Anti-nodes oscillate between these points of maximum displacement.

\subsubsection{Harmonics}
% Something i don't really understand *and* don't really understand the point of
The first harmonic is half a wavelength. It is the shortest length stationary wave between the enclosures of two nodes

$$T = \frac{1}{2\pi}\sqrt{\frac{f}{\mu}}$$ % something like that; I cannot derive and thus do not know its significance

\subsection{Refraction, diffraction \& Interference}
\textbf{Path Difference ($\Delta x$)}: The difference in length that a wave has taken to reach a given point 

\textbf{Phase Difference ($\Delta\phi$)}: The difference in phase between two given points (usually at their incidence).

Example: if wave A is at $6.56\lambda$ and wave B is at $12.31\lambda$ then:
$$\Delta x = 5.75m$$
$$\Delta\phi = (.56 - .31) \times 2\pi = \frac{\pi}{2}$$

This would mean that the waves are $\frac{\pi}{2}$ out of phase. This idea is important for the upcoming experiments.

\textbf{Coherence}: A wave is coherent when it has the same wavelength, frequency and constant phase relationships. To ensure that the waves are coherent for these upcoming experiments, monochromatic (mostly laser) light was used so that clear interference patterns could be observed. 

\subsubsection{Young's Double Slit Experiment}
% I will need to also include diagrams from this (probably best if drawn myself on Onenote)
Thomas Young disproved the corpuscle (\textit{light was a particle}) theory of light with this experiment.

A monochromatic light source would be incident to two different slits a distance $s$ away from each other (where $s$ is the slit separation). Since the light would diffract and interfere with each other, oscillating bands of light and darkness would be shown on the screen that this would be shone onto. This screen would be a distance $D$ away from the first screen. 

$D >> s$ because ... 

\textbf{Note}: If the light was not monochromatic (e.g. a light bulb) then a filter would be made: first by shining the light through a single slit (\textbf{which achieves what effect??}) and polarising the resultant light to ensure that it is monochromatic (\textbf{again, is this right?}).

\end{document}
